\section{Transport on a nonzero Hermitian shell}

The difficult residues in both applications can be described on the
same set of vectors. We give the complete algebraic argument here,
including the maps that turn its group generators into actual incidence
paths. These are elementary unitary calculations, not a claim of a new
classification theorem. The proof combines the accepted unital argument
with its explicit transport supplement.

Let $K$ be a commutative field and $\sigma$ an involutive field
automorphism, written $a\mapsto\bar a$. On column vectors in $K^2$
put
\[
 h((x,y),(u,v))=x\bar v+y\bar u,
 \qquad N(x,y)=x\bar y+\bar x y,\qquad
 J=\begin{pmatrix}0&1\\1&0\end{pmatrix}.
\]
Write $\operatorname{SU}(h)$ for the determinant-one matrices
preserving this form. For $M=\left(\begin{smallmatrix}a&b\\c&d\end{smallmatrix}\right)$
with determinant one, the equation $M^{\mathsf T}J\sigma(M)=J$
is equivalent to
\[
 \bar a=a,\quad \bar d=d,\quad \bar b=-b,\quad \bar c=-c,
 \qquad ad-bc=1.                                             \tag{H1}
\]
Indeed solving for $\sigma(M)$ gives
$\sigma(M)=J(M^{\mathsf T})^{-1}J
=\left(\begin{smallmatrix}a&-b\\-c&d\end{smallmatrix}\right)$;
the calculation is reversible. Direct expansion also verifies
$N(Mv)=N(v)$ under (H1).

\paragraph{Theorem H.1 (unique unitary transport).}
If $N(v)=N(w)=k\ne0$, there is exactly one element of
$\operatorname{SU}(h)$ carrying $v$ to $w$. For $v=(x,y)$ and
$w=(u,r)$ it is
\[
 T_{w,v}=\frac1k
 \begin{pmatrix}
 u\bar y+\bar u y & u\bar x-\bar u x\\
 r\bar y-\bar r y & r\bar x+\bar r x
 \end{pmatrix}.                                              \tag{H2}
\]

\emph{Proof.}
The columns of
\[
 B_v=\begin{pmatrix}x&-\bar x\\y&\bar y\end{pmatrix}
\]
are orthogonal and have respective norms $k,-k$. Its determinant
is $k$, so it is invertible. Thus $B_wB_v^{-1}$ has determinant
one, preserves the Hermitian form, and carries the first column of
$B_v$ to that of $B_w$. Matrix multiplication gives (H2). The
argument includes characteristic two: the two columns remain
independent because their determinant is the nonzero element $k$.

For uniqueness, (H1) implies that any such matrix taking $(x,y)$
to $(u,r)$ must take $(-\bar x,\bar y)$ to
$(-\bar u,\bar r)$. Its values on the basis formed by the columns
of $B_v$ are determined, so it is $B_wB_v^{-1}$.
\hfill$\square$

\section{Two isotropic directions suffice}

For $\bar t=-t$, the matrices
\[
 U_t=\begin{pmatrix}1&t\\0&1\end{pmatrix},\qquad
 V_t=\begin{pmatrix}1&0\\t&1\end{pmatrix}                     \tag{H3}
\]
satisfy (H1). Their displacements lie on the two isotropic coordinate
axes. If $N(v)\ne0$ and $t\ne0$, those displacements are nonzero:
neither coordinate of $v$ can be zero.

Now let $K=\mathbb F_{q^2}$ with $\sigma(a)=a^q$. The trace-zero
space has dimension one over $\mathbb F_q$. Choose a nonzero
$\delta$ with $\bar\delta=-\delta$, including in characteristic
two. Then $\delta^2\in\mathbb F_q^*$, and the entries of (H1)
satisfy $a,d\in\mathbb F_q$, $b,c\in\delta\mathbb F_q$.
Conjugation by $\operatorname{diag}(\delta,1)$ identifies these
matrices with $\operatorname{SL}(2,q)$: explicitly, conjugating
by its inverse on the left changes the off-diagonal entries to
$b/\delta$ and $c\delta$, which range over $\mathbb F_q$.
The two families (H3) therefore become all upper and lower
elementary matrices over $\mathbb F_q$.

For clarity, these elementary matrices generate $\operatorname{SL}(2,F)$
over every field $F$. If
$A=\left(\begin{smallmatrix}a&b\\c&d\end{smallmatrix}\right)$
has $a\ne0$, left multiplication by $V_{-c/a}$ and right
multiplication by $U_{-b/a}$ give $\operatorname{diag}(a,a^{-1})$.
If $a=0$, then $c\ne0$ and $U_1A$ has nonzero upper-left entry.
Finally, for $t\ne0$,
\[
 W(t)=U_tV_{-t^{-1}}U_t=
 \begin{pmatrix}0&t\\-t^{-1}&0\end{pmatrix},\qquad
 W(t)W(-1)=\operatorname{diag}(t,t^{-1}).                     \tag{H4}
\]
This proves generation, including over $\mathbb F_2$.

\paragraph{Theorem H.2 (connectivity using specified directions).}
Let $W$ be a nondegenerate Hermitian plane over $\mathbb F_{q^2}$,
let $k\in\mathbb F_q^*$, and let $\mathcal A$ be a set of its
isotropic one-spaces containing at least two distinct members.
The graph on $X_k=\{v\in W:h(v,v)=k\}$ which joins distinct
$v,w$ whenever $K(v-w)\in\mathcal A$ is connected. The group
$\operatorname{SU}(W)$ acts regularly on $X_k$, and
$|X_k|=q(q^2-1)$.

\emph{Proof.}
Choose isotropic representatives $e,f$ for two specified directions.
Their pairing is nonzero, since the orthogonal complement of an
isotropic line in a nondegenerate plane is that same line. Rescale
them to give Gram matrix $J$. Theorem H.1 proves regularity of
the special unitary action on each nonempty nonzero shell. Such
shells are nonempty: with $y=1$ the equation $N(x,1)=k$ is the
surjective trace equation. The two root families in this chosen
basis generate the whole group by (H4), and each nonidentity
root element makes a nonzero move in one of the two allowed
directions. An element carrying one shell vector to another can
therefore be written as a product of moves along graph edges.
This proves connectivity. Finally $|\operatorname{SL}(2,q)|
=q(q^2-1)$, by choosing a nonzero first column and then one of
the $q$ second columns with determinant one. Regularity gives
the claimed shell size. \hfill$\square$

\section{The projective incidence map}

Let $V$ be a nondegenerate Hermitian three-space over
$K=\mathbb F_{q^2}$. Its classical unital has as points the
isotropic one-spaces, and as blocks their intersections with
secant projective lines. It is the classical
$S(2,q+1,q^3+1)$ design, with $q^2$ blocks through a point.
These design parameters and elementary finite-field Hermitian
linear algebra are the classical inputs used here.

Fix a block $L$. Its projective line corresponds to a nondegenerate
two-space $W$. A degenerate two-space would have a one-dimensional
radical and just one isotropic projective point, whereas a secant
has $q+1$. Write $V=W\perp Kz$ and use surjectivity of the field
norm to arrange $h(z,z)=1$. Every unital point outside $L$ has a
unique representative $v+z$, where
\[
                    v\in W,\qquad h(v,v)=-1.                \tag{H5}
\]
For distinct such vectors $v,w$, the projective line joining
$\langle v+z\rangle$ and $\langle w+z\rangle$ meets $W$ in
$\langle v-w\rangle$. Hence their joining block meets $L$ at
an isotropic point precisely when $v-w$ is isotropic. The point
of intersection specifies its direction, with no arbitrary
choice of a group isomorphism.

\paragraph{Corollary H.3 (full and deleted pencils).}
Let $\mathcal A$ be a set of at least two points on the unital
block $L$. Consider the bipartite geometry whose first part is
the unital points outside $L$, whose second part consists of
flags $(a,M)$ with $a\in\mathcal A$, $M\ne L$, and whose
incidence is $Q\in M$. This graph is connected for every prime
power $q\geq2$.

\emph{Proof.}
By (H5), its point-collinearity graph is exactly the shell graph
of Theorem H.2 for the directions represented by $\mathcal A$.
Every flag vertex has $q$ outside-point neighbors, so connectivity
of the point-collinearity graph connects the entire bipartite
graph. \hfill$\square$

Taking $\mathcal A=L$ gives the block residue of the Van Nypelseer
construction. Taking $\mathcal A=L\setminus\{p\}$ gives the
triangle residue in Part I, Theorem 4. Since $|L|=q+1$, both choices
retain two directions even at $q=2$. The proof thus justifies this
particular pencil deletion; it does not assume that arbitrary
deletion preserves connectivity.

\section{Formal scope and attribution}

The complete norm-shell proof was developed earlier in this project
for the conjectural statement in [VN, Section 4.6]. Its acceptance
and the question accounting are prior work; this edition does not
count that answer again. Formula (H2), the elementary unitary
calculation and the classical root generators are not claimed to
be historically new.

The retained module \texttt{HermitianPlaneTransport.lean} proves
thirteen statements: the Gram-equation equivalence, invariance of
the norm under the involution, norm preservation, existence and
uniqueness of (H2), and the root-matrix and nonzero-displacement
identities. They hold over an arbitrary field with an involution.
The retained \texttt{ResidueConnectivity.lean} supplies the abstract
generating-move connectivity implication. The finite-field
identification with $\operatorname{SL}(2,q)$, elementary generation,
projective map (H5) and the source geometry applications are
written proofs here. They are not formal consequences claimed to
have been instantiated in those modules.
